\documentclass[aspectratio=169,11pt]{beamer}

% =========================================================
% TEMPLATE
% =========================================================
\input{~/Dropbox/Latex/TemplateSlides/template-slide.tex}

% =========================================================
% INFORMACIÓN DEL CURSO
% =========================================================
\renewcommand{\CourseName}{Programación con JavaScript}
\renewcommand{\CourseCode}{ADSO - CTMA}
\renewcommand{\Faculty}{Análisis y Desarrollo de Software}
\renewcommand{\Institution}{Servicio Nacional de Aprendizaje SENA}
\renewcommand{\Instructor}{Deimer Miranda Montoya, MSc.(c).}
\renewcommand{\Email}{damiranda@sena.edu.co}
\renewcommand{\InstitutionLogo}{Template/LogoSena.png}

% =========================================================
% COMANDOS AUXILIARES
% =========================================================
\newcommand{\concept}[1]{%
\textcolor{BrandBlue}{\textbf{#1}}%
}

\newcommand{\file}[1]{%
\texttt{#1}%
}

% =========================================================
% META
% =========================================================
\title{\CourseName}
\author{\Instructor}
\date{\today}

% =========================================================
% DOCUMENTO
% =========================================================
\begin{document}

% =========================================================
% PORTADA
% =========================================================
\TitleSlide

% =========================================================
% NORMAS
% =========================================================
\NormasSlide

% =========================================================
% REFLEXIÓN INICIAL
% =========================================================
\begin{frame}{Reflexión inicial}
    \centering

    \textbf{\large
        "Aprender algo nuevo no siempre significa empezar desde cero;\\
        muchas veces significa conectar lo que ya sabemos."}



    \smartimage[4cm]{Template/person.png}
\end{frame}

%==========================================================================================================================
\SectionSlide{Clase 4}{Condicionales, eventos e interacción dinámica con el DOM}
%==========================================================================================================================


% =========================================================
% RUTA DE LA CLASE
% =========================================================
\begin{frame}{Ruta de trabajo}

    \Idea[Partimos de conocimientos previos]{
        Ya conocemos variables, operadores, condicionales y lógica básica
        gracias a Python. En esta clase nos concentraremos en trasladar
        esos conceptos a JavaScript y utilizarlos para construir
        interfaces web interactivas.
    }



    \begin{center}
        \begin{tikzpicture}[
            node distance=0.65cm,
            every node/.style={font=\small}
            ]

            \node[
            draw,
            rounded corners,
            fill=IdeaBody,
            text width=2cm,
            align=center,
            minimum height=0.9cm
            ] (dom) {
                DOM
            };

            \node[
            draw,
            rounded corners,
            fill=DefBody,
            right=of dom,
            text width=2cm,
            align=center,
            minimum height=0.9cm
            ] (evento) {
                Evento
            };

            \node[
            draw,
            rounded corners,
            fill=QuestionBody,
            right=of evento,
            text width=2cm,
            align=center,
            minimum height=0.9cm
            ] (dato) {
                Dato
            };

            \node[
            draw,
            rounded corners,
            fill=PracticeBody,
            right=of dato,
            text width=2cm,
            align=center,
            minimum height=0.9cm
            ] (decision) {
                Condición
            };

            \node[
            draw,
            rounded corners,
            fill=ResultBody,
            right=of decision,
            text width=2cm,
            align=center,
            minimum height=0.9cm
            ] (salida) {
                Interfaz
            };

            \draw[-{Latex[length=2.5mm]}, thick] (dom) -- (evento);
            \draw[-{Latex[length=2.5mm]}, thick] (evento) -- (dato);
            \draw[-{Latex[length=2.5mm]}, thick] (dato) -- (decision);
            \draw[-{Latex[length=2.5mm]}, thick] (decision) -- (salida);

        \end{tikzpicture}
    \end{center}

\end{frame}


% =========================================================
% OBJETIVOS
% =========================================================
\begin{frame}{¿Qué vamos a construir?}

    \Result[Al finalizar la clase]{
        El aprendiz podrá:
        \begin{itemize}
            \item trasladar estructuras condicionales de Python a JavaScript;
            \item diferenciar \texttt{==} de \texttt{===};
            \item usar \texttt{addEventListener()};
            \item reconocer los eventos más usados del navegador;
            \item actualizar una interfaz mientras el usuario interactúa;
            \item validar campos mediante eventos;
            \item procesar correctamente un formulario con \texttt{submit}.
        \end{itemize}
    }

\end{frame}


%==========================================================================================================================
\SectionSlide{Bloque 1}{Decisiones: de Python a JavaScript}
%==========================================================================================================================


% =========================================================
% PYTHON VS JS
% =========================================================
\begin{frame}[fragile]{El concepto no cambia, la sintaxis sí}

    \begin{columns}[T,onlytextwidth]

        \begin{column}{0.45\textwidth}

            \DefBlock[Python]{
                La estructura depende de indentación.
            }

            \begin{consolecode}
            if temperatura > 30:
                print("Hace calor")
            else:
                print("Temperatura moderada")
            \end{consolecode}

        \end{column}

        \begin{column}{0.5\textwidth}

            \DefBlock[JavaScript]{
                La estructura utiliza paréntesis y llaves.
            }

            \begin{javascriptcode}
                if (temperatura > 30) {
                    console.log("Hace calor");
                } else {
                    console.log("Temperatura moderada");
                }
            \end{javascriptcode}

        \end{column}

    \end{columns}

\end{frame}


% =========================================================
% IF ELSE IF
% =========================================================
\begin{frame}[fragile]{Selección múltiple con \texttt{else if}}

    \DefBlock[En JavaScript]{
        Podemos encadenar condiciones mediante
        \texttt{if}, \texttt{else if} y \texttt{else}.
    }



    \begin{javascriptcode}
        let nivel = 72;

        if (nivel >= 90) {
            console.log("Nivel alto");
        } else if (nivel >= 60) {
            console.log("Nivel medio");
        } else {
            console.log("Nivel bajo");
        }
    \end{javascriptcode}

    \vspace{-0.25cm}

    \Idea[Recordatorio]{
        Las condiciones se evalúan de arriba hacia abajo.
        Cuando una condición es verdadera,
        las alternativas posteriores no se ejecutan.
    }

\end{frame}


% =========================================================
% RELACIONALES
% =========================================================
\begin{frame}{Operadores que utilizaremos}

    \begin{center}
        \begin{tabular}{
            cll
        }

        \toprule
        \textbf{Operador} & \textbf{Significado} & \textbf{Ejemplo} \\
        \midrule

        \texttt{>}
        & Mayor que
        & \texttt{edad > 18}
        \\[0.12cm]

        \texttt{>=}
        & Mayor o igual
        & \texttt{nota >= 3}
        \\[0.12cm]

        \texttt{<}
        & Menor que
        & \texttt{stock < 5}
        \\[0.12cm]

        \texttt{<=}
        & Menor o igual
        & \texttt{nivel <= 100}
        \\[0.12cm]

        \texttt{===}
        & Igualdad estricta
        & \texttt{opcion === 1}
        \\[0.12cm]

        \texttt{!==}
        & Diferencia estricta
        & \texttt{estado !== "activo"}
        \\

        \bottomrule

    \end{tabular}
    \end{center}

\end{frame}


% =========================================================
% IGUALDAD ESTRICTA
% =========================================================
\begin{frame}[fragile]{Igualdad estricta: \texttt{===}}

    \DefBlock[\texttt{===}]{
        Compara simultáneamente el valor y el tipo de dato.
        La expresión solo será verdadera si ambos coinciden.
    }



    \begin{javascriptcode}
        console.log(5 === 5);
        console.log(5 === "5");
        console.log("admin" === "admin");
        console.log(true === 1);
    \end{javascriptcode}



    \Result[Lectura]{
        \texttt{5 === "5"} produce \texttt{false},
        porque un valor es \texttt{number}
        y el otro es \texttt{string}.
    }

\end{frame}


% =========================================================
% ESTRICTA VS NO ESTRICTA
% =========================================================
\begin{frame}[fragile]{\texttt{===} frente a \texttt{==}}

    \DefBlock[Diferencia principal]{
        \texttt{==} puede convertir automáticamente los tipos
        antes de comparar. \texttt{===} no realiza esa conversión.
    }

    \vspace{-0.25cm}

    \begin{javascriptcode}
        console.log(5 == "5");
        console.log(5 === "5");
        console.log(0 == false);
        console.log(0 === false);
    \end{javascriptcode}



    \begin{center}
        \begin{tabular}{
            lll
        }

        \toprule
        \textbf{Valores} & \textbf{\texttt{==}} & \textbf{\texttt{===}} \\
        \midrule
        \texttt{5} y \texttt{"5"} & \texttt{true} & \texttt{false} \\[0.12cm]
        \texttt{0} y \texttt{false} & \texttt{true} & \texttt{false} \\
        \bottomrule

    \end{tabular}
    \end{center}



    \Idea[Buena práctica]{
        Utilizaremos preferentemente \texttt{===} y \texttt{!==}.
    }

\end{frame}


% =========================================================
% OPERADORES LÓGICOS
% =========================================================
\begin{frame}{Operadores lógicos}

    \begin{center}
        \begin{tabular}{
            c
            c
            p{0.5\linewidth}
        }

        \toprule
        \textbf{JS} & \textbf{Python} & \textbf{Uso} \\
        \midrule

        \texttt{\&\&}
        & \texttt{and}
        & Ambas condiciones deben cumplirse.
        \\[0.18cm]

        \texttt{||}
        & \texttt{or}
        & Basta que una condición se cumpla.
        \\[0.18cm]

        \texttt{!}
        & \texttt{not}
        & Niega una expresión booleana.
        \\

        \bottomrule

    \end{tabular}
    \end{center}



    \Result[Ejemplo]{
        \texttt{edad >= 18 \&\& edad <= 60}
    }

\end{frame}


% =========================================================
% EJEMPLO LÓGICOS
% =========================================================
\begin{frame}[fragile]{Ejemplo: rango de operación}

    \DefBlock[Situación]{
        Un sensor debe considerarse dentro del rango normal
        solamente entre 20 y 80 unidades.
    }



    \begin{javascriptcode}
        let lectura = 63;

        if (lectura >= 20 && lectura <= 80) {
            console.log("Lectura dentro del rango");
        } else {
            console.log("Lectura fuera del rango");
        }
    \end{javascriptcode}



    \Idea[Traducción mental]{
        \texttt{\&\&} puede leerse como ``y''.
    }

\end{frame}


%==========================================================================================================================
\SectionSlide{Bloque 2}{Eventos y addEventListener()}
%==========================================================================================================================


% =========================================================
% DE CÓDIGO LINEAL A INTERACCIÓN
% =========================================================
\begin{frame}{¿Qué cambia en una aplicación web?}

    \Idea[Programa tradicional]{
        El programa ejecuta instrucciones siguiendo un flujo previsto.
    }



    \Idea[Aplicación web]{
        Muchas acciones ocurren cuando el usuario interactúa:
        escribe, hace clic, cambia una opción, entra a un campo
        o envía un formulario.
    }



    \begin{center}
        \begin{tikzpicture}[
            node distance=0.5cm,
            every node/.style={font=\small}
            ]

            \node[
            draw,
            rounded corners,
            fill=IdeaBody,
%            text width=2.3cm,
            align=center,
            minimum height=0.9cm
            ] (usuario) {
                Usuario
            };

            \node[
            draw,
            rounded corners,
            fill=QuestionBody,
            right=of usuario,
%            text width=2.3cm,
            align=center,
            minimum height=0.9cm
            ] (evento) {
                Evento
            };

            \node[
            draw,
            rounded corners,
            fill=PracticeBody,
            right=of evento,
%            text width=2.6cm,
            align=center,
            minimum height=0.9cm
            ] (listener) {
                \texttt{addEventListener()}
            };

            \node[
            draw,
            rounded corners,
            fill=DefBody,
            right=of listener,
%            text width=2.3cm,
            align=center,
            minimum height=0.9cm
            ] (funcion) {
                Función
            };

            \node[
            draw,
            rounded corners,
            fill=ResultBody,
            right=of funcion,
%            text width=2.3cm,
            align=center,
            minimum height=0.9cm
            ] (dom) {
                DOM cambia
            };

            \draw[-{Latex[length=2.5mm]}, thick] (usuario) -- (evento);
            \draw[-{Latex[length=2.5mm]}, thick] (evento) -- (listener);
            \draw[-{Latex[length=2.5mm]}, thick] (listener) -- (funcion);
            \draw[-{Latex[length=2.5mm]}, thick] (funcion) -- (dom);

        \end{tikzpicture}
    \end{center}

\end{frame}


% =========================================================
% QUÉ ES EVENTO
% =========================================================
\begin{frame}{¿Qué es un evento?}

    \DefBlock[Evento]{
        Es una acción o suceso que ocurre en la página
        y que JavaScript puede detectar.
    }



    \ExampleBlock[Ejemplos]{
        \begin{itemize}
            \item El usuario hace clic.
            \item Escribe dentro de un campo.
            \item Selecciona una opción.
            \item Entra o sale de un campo.
            \item Presiona una tecla.
            \item Envía un formulario.
        \end{itemize}
    }

\end{frame}


% =========================================================
% SINTAXIS EVENTLISTENER
% =========================================================
\begin{frame}[fragile]{Estructura de \texttt{addEventListener()}}

    \begin{javascriptcode}
        elemento.addEventListener("evento", function () {
            // instrucciones que se ejecutarán
        });
    \end{javascriptcode}



\begin{center}
        \begin{tabular}{
            ll
        }

        \toprule
        \textbf{Parte} & \textbf{Función} \\
        \midrule

        \texttt{elemento}
        & Referencia a un elemento del DOM.
        \\[0.14cm]

        \texttt{addEventListener()}
        & Registra un escuchador de eventos.
        \\[0.14cm]

        \texttt{"evento"}
        & Indica qué acción queremos detectar.
        \\[0.14cm]

        \texttt{function()}
        & Contiene el código que debe ejecutarse.
        \\

        \bottomrule

    \end{tabular}
\end{center}

\end{frame}


% =========================================================
% EVENTOS A ESTUDIAR
% =========================================================
\begin{frame}{Eventos que utilizaremos}

\begin{center}
        \begin{tabular}{
            ll
        }

        \toprule
        \textbf{Evento} & \textbf{Cuándo ocurre} \\
        \midrule

        \texttt{click}
        & Al hacer clic sobre un elemento.
        \\[0.12cm]

        \texttt{input}
        & Cada vez que cambia el contenido de un campo.
        \\[0.12cm]

        \texttt{change}
        & Cuando se confirma un cambio de valor.
        \\[0.12cm]

        \texttt{focus}
        & Cuando un elemento recibe el foco.
        \\[0.12cm]

        \texttt{blur}
        & Cuando un elemento pierde el foco.
        \\[0.12cm]

        \texttt{keydown}
        & Cuando se presiona una tecla.
        \\[0.12cm]

        \texttt{submit}
        & Cuando se intenta enviar un formulario.
        \\

        \bottomrule

    \end{tabular}
\end{center}

\end{frame}
%==========================================================================================================================
\SectionSlide{Bloque 3}{Eventos de interacción}
%==========================================================================================================================


% =========================================================
% CLICK
% =========================================================
\begin{frame}[fragile]{Evento \texttt{click}}

    \DefBlock[\texttt{click}]{
        Se ejecuta cuando el usuario hace clic sobre el elemento.
        Es frecuente en botones.
    }



    \begin{htmlcode}
        <button id="mostrar">Mostrar mensaje</button>
        <p id="salida"></p>
    \end{htmlcode}



    \begin{javascriptcode}
        const boton = document.getElementById("mostrar");
        const salida = document.getElementById("salida");

        boton.addEventListener("click", function () {
            salida.textContent = "El botón fue presionado.";
        });
    \end{javascriptcode}

\end{frame}


% =========================================================
% INPUT
% =========================================================
\begin{frame}[fragile]{Evento \texttt{input}: actualización inmediata}

    \DefBlock[\texttt{input}]{
        Se dispara cada vez que cambia el contenido
        de un campo de entrada.
    }



    \begin{htmlcode}
        <input id="nombre" type="text">
        <p id="vista"></p>
    \end{htmlcode}



    \begin{javascriptcode}
        const nombre = document.getElementById("nombre");
        const vista = document.getElementById("vista");

        nombre.addEventListener("input", function () {
            vista.textContent = nombre.value;
        });
    \end{javascriptcode}



    \Idea[Característica]{
        No es necesario pulsar un botón.
        El resultado cambia mientras el usuario escribe.
    }

\end{frame}


% =========================================================
% FLUJO INPUT
% =========================================================
\begin{frame}{¿Por qué \texttt{input} resulta tan útil?}

    \centering

    \begin{tikzpicture}[
        node distance=0.5cm,
        every node/.style={font=\small}
        ]

        \node[
        draw,
        rounded corners,
        fill=IdeaBody,
        text width=2.3cm,
        align=center,
        minimum height=0.9cm
        ] (escribe) {
            Usuario\\
            escribe
        };

        \node[
        draw,
        rounded corners,
        fill=QuestionBody,
        right=of escribe,
        text width=2.3cm,
        align=center,
        minimum height=0.9cm
        ] (input) {
            Evento\\
            \texttt{input}
        };

        \node[
        draw,
        rounded corners,
        fill=DefBody,
        right=of input,
        text width=2.3cm,
        align=center,
        minimum height=0.9cm
        ] (value) {
            Leer\\
            \texttt{.value}
        };

        \node[
        draw,
        rounded corners,
        fill=PracticeBody,
        right=of value,
        text width=2.3cm,
        align=center,
        minimum height=0.9cm
        ] (process) {
            Procesar
        };

        \node[
        draw,
        rounded corners,
        fill=ResultBody,
        right=of process,
        text width=2.3cm,
        align=center,
        minimum height=0.9cm
        ] (dom) {
            Actualizar\\
            DOM
        };

        \draw[-{Latex[length=2.5mm]}, thick] (escribe) -- (input);
        \draw[-{Latex[length=2.5mm]}, thick] (input) -- (value);
        \draw[-{Latex[length=2.5mm]}, thick] (value) -- (process);
        \draw[-{Latex[length=2.5mm]}, thick] (process) -- (dom);

    \end{tikzpicture}



    \Result[Resultado]{
        La interfaz puede reaccionar inmediatamente
        sin esperar un clic adicional.
    }

\end{frame}


% =========================================================
% EJEMPLO INPUT CONTADOR
% =========================================================
\begin{frame}[fragile]{Ejemplo: contador de caracteres}

    \DefBlock[Situación]{
        Mostrar cuántos caracteres ha escrito el usuario.
    }



    \begin{htmlcode}
        <input id="mensaje" type="text">
        <p id="contador"></p>
    \end{htmlcode}



    \begin{javascriptcode}
        const mensaje = document.getElementById("mensaje");
        const contador = document.getElementById("contador");

        mensaje.addEventListener("input", function () {
            let cantidad = mensaje.value.length;

            contador.textContent =
            "Caracteres escritos: " + cantidad;
        });
    \end{javascriptcode}

\end{frame}


% =========================================================
% INPUT CON IF
% =========================================================
\begin{frame}[fragile]{Condicional + \texttt{input}}

    \DefBlock[Interacción dinámica]{
        Podemos evaluar el valor cada vez que el usuario lo modifica.
    }



    \begin{javascriptcode}
        const entradaNivel = document.getElementById("nivel");
        const salidaNivel = document.getElementById("estado");

        entradaNivel.addEventListener("input", function () {
            let nivel = Number(entradaNivel.value);

            if (nivel >= 80) {
                salidaNivel.textContent = "Nivel alto";
            } else if (nivel >= 40) {
                salidaNivel.textContent = "Nivel medio";
            } else {
                salidaNivel.textContent = "Nivel bajo";
            }
        });
    \end{javascriptcode}

\end{frame}


% =========================================================
% MICROPRÁCTICA 1
% =========================================================
\begin{frame}{Micropráctica 1}

    \Practice[Indicador de aforo]{
        Crear un campo que permita ingresar la cantidad
        de personas presentes en un lugar.
        La salida debe actualizarse mientras se escribe.
    }



    \Result[Clasificación]{
        \begin{itemize}
            \item menos de 20 personas: \texttt{Aforo bajo};
            \item entre 20 y 49: \texttt{Aforo medio};
            \item 50 o más: \texttt{Aforo alto}.
        \end{itemize}
    }



    \Idea[Condición]{
        No utilizar un botón: emplear el evento \texttt{input}.
    }

\end{frame}


% =========================================================
% CHANGE
% =========================================================
\begin{frame}[fragile]{Evento \texttt{change}}

    \DefBlock[\texttt{change}]{
        Se utiliza cuando queremos reaccionar
        a un cambio confirmado en un control.
        Es especialmente útil con \texttt{select}.
    }



    \begin{htmlcode}
    <select id="transporte">
        <option value="">Seleccione</option>
        <option value="bus">Bus</option>
        <option value="metro">Metro</option>
        <option value="bici">Bicicleta</option>
    </select>
    <p id="respuesta"></p>
    \end{htmlcode}

    \vspace{0.1cm}

    \begin{javascriptcode}
        transporte.addEventListener("change", function () {
            respuesta.textContent =
            "Opción seleccionada: " + transporte.value;
        });
    \end{javascriptcode}

\end{frame}


% =========================================================
% INPUT VS CHANGE
% =========================================================
\begin{frame}{\texttt{input} frente a \texttt{change}}

    \begin{tabular}{
           lll
        }

        \toprule
        & \textbf{\texttt{input}} & \textbf{\texttt{change}} \\
        \midrule

        Momento
        &
        Mientras cambia el valor.
        &
        Cuando el cambio queda establecido.
        \\[0.18cm]

        Uso frecuente
        &
        Campos de texto y número.
        &
        \texttt{select}, checkbox y otros controles.
        \\[0.18cm]

        Sensación
        &
        Tiempo real.
        &
        Cambio confirmado.
        \\

        \bottomrule

    \end{tabular}



    \Idea[No son enemigos]{
        Elegimos el evento según el comportamiento
        que queremos para la interfaz.
    }

\end{frame}


% =========================================================
% CHANGE CON CONDICIONAL
% =========================================================
\begin{frame}[fragile]{Ejemplo: método de entrega}

    \DefBlock[Situación]{
        Mostrar información diferente según
        la opción elegida por el usuario.
    }

%<label for="tipoEntrega">Tipo de entrega:</label>
%
%<select id="tipoEntrega">
%<option value="">Seleccione una opción</option>
%<option value="domicilio">Entrega a domicilio</option>
%<option value="tienda">Recoger en tienda</option>
%</select>
%
%<p id="salida">Seleccione una opción.</p>


    \begin{javascriptcode}
        tipoEntrega.addEventListener("change", function () {
            let opcion = tipoEntrega.value;

            if (opcion === "domicilio") {
                salida.textContent = "Entrega estimada: 2 días.";
            } else if (opcion === "tienda") {
                salida.textContent = "Puede recoger hoy.";
            } else {
                salida.textContent = "Seleccione una opción.";
            }
        });
    \end{javascriptcode}

\end{frame}


% =========================================================
% MICROPRÁCTICA 2
% =========================================================
\begin{frame}{Micropráctica 2}

    \Practice[Selector de prioridad]{
        Crear un \texttt{select} con tres opciones:
        baja, media y alta.
    }



    \Result[Comportamiento]{
        Al cambiar la selección:
        \begin{itemize}
            \item baja: mostrar \texttt{Puede esperar};
            \item media: mostrar \texttt{Atender próximamente};
            \item alta: mostrar \texttt{Atención prioritaria}.
        \end{itemize}
    }



    \Idea[Evento]{
        Utilizar \texttt{change}.
    }

\end{frame}


% =========================================================
% FOCUS
% =========================================================
\begin{frame}[fragile]{Evento \texttt{focus}}

    \DefBlock[\texttt{focus}]{
        Ocurre cuando un elemento recibe el foco;
        por ejemplo, cuando el usuario hace clic dentro de un campo.
    }



    \begin{htmlcode}
        <input id="usuario" type="text">
        <p id="ayuda"></p>
    \end{htmlcode}



    \begin{javascriptcode}
        usuario.addEventListener("focus", function () {
            ayuda.textContent =
            "Use entre 5 y 12 caracteres.";
        });
    \end{javascriptcode}



    \Idea[Aplicación]{
        Es útil para mostrar ayuda contextual.
    }

\end{frame}


% =========================================================
% BLUR
% =========================================================
\begin{frame}[fragile]{Evento \texttt{blur}}

    \DefBlock[\texttt{blur}]{
        Ocurre cuando un elemento pierde el foco.
        Es útil para revisar un campo después de que el usuario termina.
    }



    \begin{javascriptcode}
        usuario.addEventListener("blur", function () {
            let cantidad = usuario.value.length;

            if (cantidad < 5) {
                ayuda.textContent = "Nombre demasiado corto.";
            } else {
                ayuda.textContent = "Nombre válido.";
            }
        });
    \end{javascriptcode}

\end{frame}


% =========================================================
% FOCUS BLUR DIAGRAM
% =========================================================
\begin{frame}{Ciclo \texttt{focus} -- \texttt{blur}}

    \centering

    \begin{tikzpicture}[
        node distance=1.1cm,
        every node/.style={font=\small}
        ]

        \node[
        draw,
        rounded corners,
        fill=IdeaBody,
        text width=2.7cm,
        align=center,
        minimum height=0.9cm
        ] (fuera) {
            Campo sin foco
        };

        \node[
        draw,
        rounded corners,
        fill=PracticeBody,
        right=of fuera,
        text width=2.7cm,
        align=center,
        minimum height=0.9cm
        ] (focus) {
            \texttt{focus}\\
            mostrar ayuda
        };

        \node[
        draw,
        rounded corners,
        fill=DefBody,
        right=of focus,
        text width=2.7cm,
        align=center,
        minimum height=0.9cm
        ] (editar) {
            Usuario escribe
        };

        \node[
        draw,
        rounded corners,
        fill=ResultBody,
        right=of editar,
        text width=2.7cm,
        align=center,
        minimum height=0.9cm
        ] (blur) {
            \texttt{blur}\\
            validar
        };

        \draw[-{Latex[length=2.5mm]}, thick] (fuera) -- (focus);
        \draw[-{Latex[length=2.5mm]}, thick] (focus) -- (editar);
        \draw[-{Latex[length=2.5mm]}, thick] (editar) -- (blur);

    \end{tikzpicture}



    \Idea[Patrón útil]{
        \texttt{focus}: orientar al usuario.
        \qquad
        \texttt{blur}: validar lo que escribió.
    }

\end{frame}
%==========================================================================================================================
\SectionSlide{Bloque 4}{Teclado, formularios y objeto event}
%==========================================================================================================================


% =========================================================
% KEYDOWN
% =========================================================
\begin{frame}[fragile]{Evento \texttt{keydown}}

    \DefBlock[\texttt{keydown}]{
        Se dispara cuando el usuario presiona una tecla
        mientras un elemento tiene el foco.
    }



    \begin{javascriptcode}
        entrada.addEventListener("keydown", function (event) {
            console.log(event.key);
        });
    \end{javascriptcode}



    \Idea[Nuevo elemento]{
        La función puede recibir un objeto denominado \texttt{event}
        con información sobre lo que ocurrió.
    }

\end{frame}


% =========================================================
% EVENT OBJECT
% =========================================================
\begin{frame}{El objeto \texttt{event}}

    \DefBlock[\texttt{event}]{
        Representa el evento que acaba de ocurrir
        y contiene información relacionada con él.
    }



\begin{center}
        \begin{tabular}{
            l
            p{0.52\linewidth}
        }

        \toprule
        \textbf{Elemento} & \textbf{Información} \\
        \midrule

        \texttt{event.key}
        &
        Tecla presionada.
        \\[0.15cm]

        \texttt{event.target}
        &
        Elemento que produjo el evento.
        \\[0.15cm]

        \texttt{event.type}
        &
        Tipo de evento ocurrido.
        \\[0.15cm]

        \texttt{event.preventDefault()}
        &
        Evita el comportamiento predeterminado del navegador.
        \\

        \bottomrule

    \end{tabular}
\end{center}

\end{frame}


% =========================================================
% ENTER
% =========================================================
\begin{frame}[fragile]{Detectar una tecla específica}

    \DefBlock[Ejemplo]{
        Podemos ejecutar una acción únicamente
        cuando la tecla presionada sea \texttt{Enter}.
    }



    \begin{javascriptcode}
        busqueda.addEventListener("keydown", function (event) {

            if (event.key === "Enter") {
                salida.textContent =
                "Buscando: " + busqueda.value;
            }

        });
    \end{javascriptcode}



    \Result[Combinación]{
        Evento de teclado + igualdad estricta + condicional + DOM.
    }

\end{frame}


% =========================================================
% MICROPRÁCTICA 3
% =========================================================
\begin{frame}{Micropráctica 3}

    \Practice[Campo de búsqueda]{
        Crear un campo de texto donde el usuario escriba
        una palabra.
    }



    \Result[Comportamiento]{
        \begin{itemize}
            \item Mientras escribe, mostrar el texto actual.
            \item Si presiona \texttt{Enter}, mostrar:
            \texttt{Buscando: ...}
        \end{itemize}
    }



    \Idea[Eventos]{
        Utilizar \texttt{input} y \texttt{keydown}.
    }

\end{frame}


% =========================================================
% FORMULARIOS
% =========================================================
\begin{frame}{¿Por qué trabajar con \texttt{submit}?}

    \Idea[Un error frecuente]{
        Cuando existe un formulario, no conviene depender
        únicamente del clic sobre el botón.
    }



    \DefBlock[Formulario]{
        El evento \texttt{submit} representa la intención
        de enviar el formulario completo.
    }



    \Result[Ventaja]{
        El formulario puede enviarse mediante un botón
        o mediante acciones equivalentes del navegador,
        como presionar \texttt{Enter}.
    }

\end{frame}


% =========================================================
% HTML FORM
% =========================================================
\begin{frame}[fragile]{Estructura básica de un formulario}

    \begin{htmlcode}
    <form id="formulario">

        <label for="correo">Correo:</label>
        <input id="correo" type="email">

        <button type="submit">
        Enviar
        </button>

    </form>

    <p id="mensaje"></p>
    \end{htmlcode}



    \Idea[Elemento importante]{
        El botón utiliza \texttt{type="submit"}
        porque pertenece al formulario.
    }

\end{frame}


% =========================================================
% SUBMIT
% =========================================================
\begin{frame}[fragile]{Evento \texttt{submit}}

    \begin{javascriptcode}
        const formulario = document.getElementById("formulario");

        formulario.addEventListener("submit", function (event) {

            event.preventDefault();

            console.log("Formulario procesado");

        });
    \end{javascriptcode}



    \DefBlock[\texttt{preventDefault()}]{
        Evita temporalmente el comportamiento predeterminado
        del navegador, como recargar la página al enviar el formulario.
    }

\end{frame}


% =========================================================
% POR QUÉ PREVENT DEFAULT
% =========================================================
\begin{frame}{Flujo de un formulario con JavaScript}

    \centering

    \begin{tikzpicture}[
        node distance=0.75cm,
        every node/.style={font=\small}
        ]

        \node[
        draw,
        rounded corners,
        fill=IdeaBody,
%        text width=2.2cm,
        align=center,
        minimum height=0.9cm
        ] (form) {
            Formulario
        };

        \node[
        draw,
        rounded corners,
        fill=QuestionBody,
        right=of form,
%        text width=2.2cm,
        align=center,
        minimum height=0.9cm
        ] (submit) {
            \texttt{submit}
        };

        \node[
        draw,
        rounded corners,
        fill=WarningBody,
        right=of submit,
%        text width=2.6cm,
        align=center,
        minimum height=0.9cm
        ] (prevent) {
            \texttt{preventDefault()}
        };

        \node[
        draw,
        rounded corners,
        fill=PracticeBody,
        right=of prevent,
%        text width=2.2cm,
        align=center,
        minimum height=0.9cm
        ] (validate) {
            Validar
        };

        \node[
        draw,
        rounded corners,
        fill=ResultBody,
        right=of validate,
%        text width=2.2cm,
        align=center,
        minimum height=0.9cm
        ] (result) {
            Responder
        };

        \draw[-{Latex[length=2.5mm]}, thick] (form) -- (submit);
        \draw[-{Latex[length=2.5mm]}, thick] (submit) -- (prevent);
        \draw[-{Latex[length=2.5mm]}, thick] (prevent) -- (validate);
        \draw[-{Latex[length=2.5mm]}, thick] (validate) -- (result);

    \end{tikzpicture}

\end{frame}


% =========================================================
% VALIDACION EN SUBMIT
% =========================================================
\begin{frame}[fragile]{Validar antes de procesar}

    \DefBlock[Ejemplo]{
        Verificar que un código tenga al menos cuatro caracteres.
    }



    \begin{javascriptcode}
        formulario.addEventListener("submit", function (event) {

            event.preventDefault();

            let codigo = entradaCodigo.value;

            if (codigo.length >= 4) {
                mensaje.textContent = "Código aceptado.";
            } else {
                mensaje.textContent =
                "El código debe tener mínimo 4 caracteres.";
            }

        });
    \end{javascriptcode}

\end{frame}


%==========================================================================================================================
\SectionSlide{Bloque 5}{Condiciones y eventos trabajando juntos}
%==========================================================================================================================


% =========================================================
% REALTIME VALIDATION
% =========================================================
\begin{frame}[fragile]{Validación en tiempo real}

    \DefBlock[Ejemplo]{
        Verificar un porcentaje mientras el usuario modifica el campo.
    }



    \begin{javascriptcode}
        porcentaje.addEventListener("input", function () {

            let valor = Number(porcentaje.value);
            if (valor < 0 || valor > 100) {
                estado.textContent = "Valor fuera de rango";
            } else {
                estado.textContent = "Valor válido";
            }

        });
    \end{javascriptcode}

\vspace{-0.25cm}

    \Result[Resultado]{
        El usuario recibe retroalimentación sin pulsar un botón.
    }

\end{frame}


% =========================================================
% EMPTY FIELD
% =========================================================
\begin{frame}[fragile]{¿Qué ocurre si el campo está vacío?}

    \DefBlock[Validación de texto]{
        Podemos comparar directamente el contenido del campo
        con una cadena vacía.
    }



    \begin{javascriptcode}
        nombre.addEventListener("blur", function () {

            if (nombre.value === "") {
                mensaje.textContent =
                "Este campo es obligatorio.";
            } else {
                mensaje.textContent =
                "Campo diligenciado.";
            }

        });
    \end{javascriptcode}

\end{frame}


% =========================================================
% COMBINACIÓN VARIOS EVENTOS
% =========================================================
\begin{frame}{Un mismo elemento puede tener varios eventos}

    \Idea[Ejemplo conceptual]{
        Un campo de correo podría:
        \begin{itemize}
            \item mostrar ayuda con \texttt{focus};
            \item actualizar información con \texttt{input};
            \item validarse con \texttt{blur};
            \item detectar \texttt{Enter} con \texttt{keydown}.
        \end{itemize}
    }



    \begin{center}
        \begin{tikzpicture}[
            node distance=0.7cm,
            every node/.style={font=\small}
            ]

            \node[
            draw,
            rounded corners,
            fill=DefBody,
            text width=2cm,
            align=center,
            minimum height=0.9cm
            ] (campo) {
                Campo
            };

            \node[
            draw,
            rounded corners,
            fill=IdeaBody,
            above right=0.55cm and 1cm of campo,
            text width=1.8cm,
            align=center
            ] (focus) {
                focus
            };

            \node[
            draw,
            rounded corners,
            fill=PracticeBody,
            below right=-0.65cm and 1cm of campo,
            text width=1.8cm,
            align=center
            ] (input) {
                input
            };

            \node[
            draw,
            rounded corners,
            fill=QuestionBody,
            below right=0.55cm and 1cm of campo,
            text width=1.8cm,
            align=center
            ] (blur) {
                blur
            };

            \node[
            draw,
            rounded corners,
            fill=ResultBody,
            right=4.1cm of campo,
            text width=2cm,
            align=center
            ] (interface) {
                Interfaz\\dinámica
            };

            \draw[-{Latex[length=2mm]}] (campo) -- (focus);
            \draw[-{Latex[length=2mm]}] (campo) -- (input);
            \draw[-{Latex[length=2mm]}] (campo) -- (blur);

            \draw[-{Latex[length=2mm]}] (focus) -- (interface);
            \draw[-{Latex[length=2mm]}] (input) -- (interface);
            \draw[-{Latex[length=2mm]}] (blur) -- (interface);

        \end{tikzpicture}
    \end{center}

\end{frame}


% =========================================================
% MICROPRÁCTICA 4
% =========================================================
\begin{frame}{Micropráctica 4}

    \Practice[Control de contraseña]{
        Crear un campo para una contraseña.
    }



    \Result[Comportamiento]{
        \begin{itemize}
            \item Con \texttt{focus}: mostrar los requisitos.
            \item Con \texttt{input}: mostrar cantidad de caracteres.
            \item Con \texttt{blur}: indicar si tiene mínimo 8 caracteres.
        \end{itemize}
    }



    \Idea[Objetivo]{
        Utilizar varios eventos sobre el mismo elemento.
    }

\end{frame}


% =========================================================
% TABLA RESUMEN
% =========================================================
\begin{frame}{¿Qué evento elegir?}

    \begin{center}
        \begin{tabular}{
            lc
        }

        \toprule
        \textbf{Necesito...} & \textbf{Evento recomendado} \\
        \midrule

        Responder a un botón
        & \texttt{click}
        \\[0.12cm]

        Actualizar mientras escribe
        & \texttt{input}
        \\[0.12cm]

        Detectar una selección
        & \texttt{change}
        \\[0.12cm]

        Mostrar ayuda al entrar
        & \texttt{focus}
        \\[0.12cm]

        Validar al salir
        & \texttt{blur}
        \\[0.12cm]

        Detectar teclas
        & \texttt{keydown}
        \\[0.12cm]

        Procesar un formulario
        & \texttt{submit}
        \\

        \bottomrule

    \end{tabular}
    \end{center}

\end{frame}
%==========================================================================================================================
\SectionSlide{Laboratorio guiado}{Registro interactivo a una actividad}
%==========================================================================================================================


% =========================================================
% LAB INTRO
% =========================================================
\begin{frame}{Objetivo del laboratorio}

    \Practice[Aplicación]{
        Construiremos un formulario de registro para una actividad.
        La interfaz reaccionará mientras el usuario interactúa
        con los diferentes campos.
    }



    \Result[Datos]{
        \begin{itemize}
            \item Nombre.
            \item Edad.
            \item Modalidad.
            \item Correo electrónico.
            \item Confirmación final mediante formulario.
        \end{itemize}
    }

\end{frame}


% =========================================================
% EVENTOS LAB
% =========================================================
\begin{frame}{Eventos utilizados en el laboratorio}

    \begin{center}
        \begin{tabular}{
            lcl
            }

            \toprule
            \textbf{Elemento} & \textbf{Evento} & \textbf{Función} \\
            \midrule

            Nombre
            & \texttt{input}
            & Vista previa en tiempo real.
            \\[0.15cm]

            Edad
            & \texttt{input}
            & Validación inmediata.
            \\[0.15cm]

            Modalidad
            & \texttt{change}
            & Mostrar información específica.
            \\[0.15cm]

            Correo
            & \texttt{focus / blur}
            & Ayuda y validación.
            \\[0.15cm]

            Formulario
            & \texttt{submit}
            & Procesar el registro completo.
            \\

            \bottomrule

        \end{tabular}
\end{center}
    \end{frame}



    % =========================================================
    % ESTRUCTURA
    % =========================================================
    \begin{frame}[fragile]{Paso 1: estructura del proyecto}

        \begin{consolecode}
            clase04-eventos/
            ├── index.html
            └── js/
                └── app.js
        \end{consolecode}



        \Idea[Separación]{
            HTML define la interfaz.
            JavaScript gestionará interacción,
            validaciones y actualización del DOM.
        }

    \end{frame}


    % =========================================================
    % HTML 1
    % =========================================================
    \begin{frame}[fragile]{Paso 2: formulario HTML}

        \begin{htmlcode}
        <form id="registro">

            <label for="nombre">Nombre:</label>
            <input id="nombre" type="text">
            <p id="vistaNombre"></p>

            <label for="edad">Edad:</label>
            <input id="edad" type="number">
            <p id="estadoEdad"></p>

            <label for="modalidad">Modalidad:</label>
            <select id="modalidad">
                <option value="">Seleccione</option>
                <option value="presencial">Presencial</option>
                <option value="virtual">Virtual</option>
            </select>

            <p id="infoModalidad"></p>
        \end{htmlcode}

    \end{frame}


    % =========================================================
    % HTML 2
    % =========================================================
    \begin{frame}[fragile]{Paso 2: completar formulario HTML}

        \begin{htmlcode}
            <label for="correo">Correo:</label>
            <input id="correo" type="email">

            <p id="ayudaCorreo"></p>

            <button type="submit">
            Registrar
            </button>

        </form>

        <p id="resultado"></p>

        <script src="js/app.js"></script>
        \end{htmlcode}

    \end{frame}


    % =========================================================
    % JS SELECTORES
    % =========================================================
    \begin{frame}[fragile]{Paso 3: referencias al DOM}

        \begin{javascriptcode}
            const formulario = document.getElementById("registro");
            const nombre = document.getElementById("nombre");
            const edad = document.getElementById("edad");
            const modalidad = document.getElementById("modalidad");
            const correo = document.getElementById("correo");

            const vistaNombre = document.getElementById("vistaNombre");
            const estadoEdad = document.getElementById("estadoEdad");
            const infoModalidad = document.getElementById("infoModalidad");
            const ayudaCorreo = document.getElementById("ayudaCorreo");
            const resultado = document.getElementById("resultado");
        \end{javascriptcode}

    \end{frame}


    % =========================================================
    % NOMBRE INPUT
    % =========================================================
    \begin{frame}[fragile]{Paso 4: nombre en tiempo real}

        \DefBlock[Evento \texttt{input}]{
            La vista previa cambia mientras el usuario escribe.
        }



        \begin{javascriptcode}
            nombre.addEventListener("input", function () {

                if (nombre.value === "") {
                    vistaNombre.textContent = "";
                } else {
                    vistaNombre.textContent =
                    "Participante: " + nombre.value;
                }

            });
        \end{javascriptcode}

    \end{frame}


    % =========================================================
    % EDAD INPUT
    % =========================================================
    \begin{frame}[fragile]{Paso 5: validar edad en tiempo real}

        \begin{javascriptcode}
            edad.addEventListener("input", function () {

                let valorEdad = Number(edad.value);

                if (valorEdad < 14 || valorEdad > 100) {
                    estadoEdad.textContent =
                    "Edad fuera del rango permitido.";
                } else {
                    estadoEdad.textContent =
                    "Edad válida.";
                }

            });
        \end{javascriptcode}



        \Idea[Combinación]{
            Aquí utilizamos DOM + \texttt{input}
            + conversión + condicional + operador lógico.
        }

    \end{frame}


    % =========================================================
    % MODALIDAD CHANGE
    % =========================================================
    \begin{frame}[fragile]{Paso 6: reaccionar a la modalidad}

        \begin{javascriptcode}
            modalidad.addEventListener("change", function () {

                if (modalidad.value === "presencial") {
                    infoModalidad.textContent = "Debe asistir a la sede asignada.";

                } else if (modalidad.value === "virtual") {
                    infoModalidad.textContent = "Recibirá un enlace de conexión.";

                } else {
                    infoModalidad.textContent = "";
                }

            });
        \end{javascriptcode}

    \end{frame}


    % =========================================================
    % CORREO FOCUS
    % =========================================================
    \begin{frame}[fragile]{Paso 7: ayuda mediante \texttt{focus}}

        \begin{javascriptcode}
            correo.addEventListener("focus", function () {

                ayudaCorreo.textContent = "Ingrese un correo electrónico válido.";

            });
        \end{javascriptcode}



        \Result[Comportamiento]{
            El mensaje aparece cuando el usuario
            comienza a trabajar con ese campo.
        }

    \end{frame}


    % =========================================================
    % CORREO BLUR
    % =========================================================
    \begin{frame}[fragile]{Paso 8: revisar al perder el foco}

        \begin{javascriptcode}
            correo.addEventListener("blur", function () {

                if (correo.value === "") {
                    ayudaCorreo.textContent = "El correo es obligatorio.";

                } else {
                    ayudaCorreo.textContent = "Correo registrado.";
                }

            });
        \end{javascriptcode}

    \end{frame}


    % =========================================================
    % SUBMIT
    % =========================================================
    \begin{frame}[fragile]{Paso 9: procesar el formulario}

        \begin{javascriptcode}
            formulario.addEventListener("submit", function (event) {

                event.preventDefault();

                let valorEdad = Number(edad.value);

                let nombreValido = nombre.value !== "";
                let edadValida = valorEdad >= 14 && valorEdad <= 100;

                let modalidadValida = modalidad.value !== "";

                let correoValido = correo.value !== "";
            \end{javascriptcode}

        \end{frame}


        % =========================================================
        % SUBMIT 2
        % =========================================================
        \begin{frame}[fragile]{Paso 10: decisión final}

            \begin{javascriptcode}
                if ( nombreValido &&
                edadValida &&
                modalidadValida &&
                correoValido
                ) {
                    resultado.textContent = "Registro completado correctamente.";

                } else {
                    resultado.textContent = "Revise los datos del formulario.";

                }

            });
        \end{javascriptcode}



        \Result[Integración]{
            El formulario completo se valida antes
            de aceptar el registro.
        }

    \end{frame}


    % =========================================================
    % DIAGRAMA LAB
    % =========================================================
    \begin{frame}{Arquitectura del laboratorio}

        \centering

        \begin{tikzpicture}[
            node distance=0.65cm,
            every node/.style={font=\small}
            ]

            \node[
            draw,
            rounded corners,
            fill=IdeaBody,
            text width=2cm,
            align=center,
            minimum height=0.9cm
            ] (html) {
                HTML
            };

            \node[
            draw,
            rounded corners,
            fill=DefBody,
            right=of html,
            text width=2cm,
            align=center,
            minimum height=0.9cm
            ] (dom) {
                DOM
            };

            \node[
            draw,
            rounded corners,
            fill=QuestionBody,
            right=of dom,
            text width=2.2cm,
            align=center,
            minimum height=0.9cm
            ] (events) {
                Eventos
            };

            \node[
            draw,
            rounded corners,
            fill=PracticeBody,
            right=of events,
            text width=2.2cm,
            align=center,
            minimum height=0.9cm
            ] (condition) {
                Condicionales
            };

            \node[
            draw,
            rounded corners,
            fill=ResultBody,
            right=of condition,
            text width=2.2cm,
            align=center,
            minimum height=0.9cm
            ] (result) {
                Respuesta
            };

            \draw[-{Latex[length=2.5mm]}, thick] (html) -- (dom);
            \draw[-{Latex[length=2.5mm]}, thick] (dom) -- (events);
            \draw[-{Latex[length=2.5mm]}, thick] (events) -- (condition);
            \draw[-{Latex[length=2.5mm]}, thick] (condition) -- (result);

        \end{tikzpicture}



        \Idea[Idea central]{
            La lógica ya conocida comienza a controlar
            el comportamiento de una aplicación web.
        }

    \end{frame}


    % =========================================================
    % RETO
    % =========================================================
    \begin{frame}{Reto de ampliación}

        \Practice[Agregar un campo de experiencia]{
            Añadir un \texttt{select} con las opciones:
            \begin{itemize}
                \item Principiante.
                \item Intermedio.
                \item Avanzado.
            \end{itemize}
        }



        \Result[Comportamiento]{
            Utilizar \texttt{change} para mostrar
            una recomendación distinta para cada nivel.
            El campo también debe formar parte
            de la validación final del formulario.
        }

    \end{frame}


    % =========================================================
    % RETO EXTRA
    % =========================================================
    \begin{frame}{Reto adicional}

        \Practice[Contador del nombre]{
            Utilizar el mismo evento \texttt{input}
            del campo nombre para mostrar también
            la cantidad de caracteres escritos.
        }



        \Result[Ejemplo esperado]{
            \texttt{Participante: Ana | 3 caracteres}
        }



        \Idea[Pista]{
            Recuerda la propiedad \texttt{.length}.
        }

    \end{frame}


    %==========================================================================================================================
    \SectionSlide{Cierre}{Síntesis y evidencia}
    %==========================================================================================================================


    % =========================================================
    % ERRORES FRECUENTES
    % =========================================================
    \begin{frame}{Errores frecuentes}

        \Warning[Revisar]{
            \begin{itemize}
                \item Leer \texttt{.value} fuera del evento cuando necesitamos el valor actualizado.
                \item Confundir \texttt{==} con \texttt{===}.
                \item Olvidar convertir campos numéricos con \texttt{Number()}.
                \item Escribir mal el nombre del evento.
                \item Colocar \texttt{submit} sobre el botón en lugar del formulario.
                \item Olvidar \texttt{event.preventDefault()}.
                \item Utilizar un \texttt{id} diferente en HTML y JavaScript.
            \end{itemize}
        }

    \end{frame}


    % =========================================================
    % EVIDENCIA
    % =========================================================
    \begin{frame}{Evidencia de la clase}

        \Practice[Producto]{
            Entregar el proyecto \file{clase04-eventos}
            con el formulario interactivo funcional.
        }



        \Result[Debe evidenciar]{
            \begin{itemize}
                \item Selección de elementos del DOM.
                \item Uso de \texttt{input}.
                \item Uso de \texttt{change}.
                \item Uso de \texttt{focus} y \texttt{blur}.
                \item Procesamiento mediante \texttt{submit}.
                \item Uso de \texttt{event.preventDefault()}.
                \item Condicionales y operadores lógicos.
                \item Actualización mediante \texttt{textContent}.
            \end{itemize}
        }

    \end{frame}


    % =========================================================
    % LISTA VERIFICACION
    % =========================================================
    \begin{frame}{Lista de verificación}

        \begin{itemize}
            \item El archivo JavaScript está enlazado correctamente.
            \item Los elementos HTML tienen identificadores únicos.
            \item Los eventos se ejecutan en el momento esperado.
            \item Los campos numéricos son convertidos con \texttt{Number()}.
            \item Se utiliza igualdad estricta cuando corresponde.
            \item Las validaciones responden mientras el usuario interactúa.
            \item El formulario no recarga la página al procesarse.
            \item El mensaje final cambia según la validez de los datos.
            \item No aparecen errores en la consola.
        \end{itemize}

    \end{frame}


    % =========================================================
    % SINTESIS EVENTOS
    % =========================================================
    \begin{frame}{Síntesis de eventos}

\begin{center}
            \begin{tabular}{
                ll
            }

            \toprule
            \textbf{Evento} & \textbf{Uso típico} \\
            \midrule

            \texttt{click} & Ejecutar una acción puntual. \\[0.10cm]
            \texttt{input} & Actualización en tiempo real. \\[0.10cm]
            \texttt{change} & Reaccionar a una selección. \\[0.10cm]
            \texttt{focus} & Mostrar ayuda al ingresar a un campo. \\[0.10cm]
            \texttt{blur} & Validar al abandonar un campo. \\[0.10cm]
            \texttt{keydown} & Detectar teclas. \\[0.10cm]
            \texttt{submit} & Procesar formularios. \\

            \bottomrule

        \end{tabular}
\end{center}

    \end{frame}


    % =========================================================
    % MAPA FINAL
    % =========================================================
    \begin{frame}{Lo que ahora podemos hacer}

        \centering

        \begin{tikzpicture}[
            node distance=0.7cm,
            every node/.style={font=\small}
            ]

            \node[
            draw,
            rounded corners,
            fill=IdeaBody,
            text width=2cm,
            align=center,
            minimum height=0.9cm
            ] (user) {
                Usuario
            };

            \node[
            draw,
            rounded corners,
            fill=QuestionBody,
            right=of user,
            text width=2cm,
            align=center,
            minimum height=0.9cm
            ] (event) {
                Evento
            };

            \node[
            draw,
            rounded corners,
            fill=DefBody,
            right=of event,
            text width=2cm,
            align=center,
            minimum height=0.9cm
            ] (value) {
                Valor
            };

            \node[
            draw,
            rounded corners,
            fill=PracticeBody,
            right=of value,
            text width=2cm,
            align=center,
            minimum height=0.9cm
            ] (if) {
                Decisión
            };

            \node[
            draw,
            rounded corners,
            fill=ResultBody,
            right=of if,
            text width=2cm,
            align=center,
            minimum height=0.9cm
            ] (dom) {
                DOM
            };

            \draw[-{Latex[length=2.5mm]}, thick] (user) -- (event);
            \draw[-{Latex[length=2.5mm]}, thick] (event) -- (value);
            \draw[-{Latex[length=2.5mm]}, thick] (value) -- (if);
            \draw[-{Latex[length=2.5mm]}, thick] (if) -- (dom);

        \end{tikzpicture}



        \Result[Progreso]{
            JavaScript ya no se limita a ejecutar cálculos:
            ahora responde a las acciones del usuario
            y modifica la interfaz dinámicamente.
        }

    \end{frame}


    % =========================================================
    % PROXIMA CLASE
    % =========================================================
    \begin{frame}{Continuación de la ruta}

        \DefBlock[Clase 5]{
            El siguiente paso será automatizar tareas
            mediante estructuras repetitivas y selección múltiple.
        }



        \Practice[Próximos contenidos]{
            \begin{itemize}
                \item \texttt{switch}.
                \item Ciclo \texttt{while}.
                \item Ciclo \texttt{do...while}.
                \item Ciclo \texttt{for}.
                \item Contadores.
                \item Acumuladores.
                \item Generación repetitiva de contenido.
            \end{itemize}
        }

    \end{frame}


    % =========================================================
    % FIN
    % =========================================================
    \begin{frame}
        \centering
        \smartimage[0.8\linewidth]{Template/gracias.png}
    \end{frame}

\end{document}